% Written by Anders Sjoqvist and Ulf Lundstrom, 2009
% The main sources are: tinyKACTL, Beta and Wikipedia
\chapter{Mathematics}

\kactlimport{basis.h}
\kactlimport{Matrix.h}

\section{Trigonometry}
\vspace{-3mm}
\begin{align*}
\sin(v+w)&{}=\sin v\cos w+\cos v\sin w\\
\cos(v+w)&{}=\cos v\cos w-\sin v\sin w\\
\tan(v+w)&{}=\dfrac{\tan v+\tan w}{1-\tan v\tan w}\\
\sin v+\sin w&{}=2\sin\dfrac{v+w}{2}\cos\dfrac{v-w}{2}\\
\cos v+\cos w&{}=2\cos\dfrac{v+w}{2}\cos\dfrac{v-w}{2}
\end{align*}
\[ (V+W)\tan(\frac{v-w}{2}){}=(V-W)\tan(\frac{v+w}{2}) \]
$V, W$ are sides opposite to angles $v, w$.
	$a\cos x+b\sin x=r\cos(x-\phi)$\\
	$a\sin x+b\cos x=r\sin(x+\phi)$\\
where $r=\sqrt{a^2+b^2}, \phi=\operatorname{atan2}(b,a)$.

\section{Geometry}

\subsection{Triangles}
Side lengths: $a,b,c$\\
Semiperimeter: $p=\dfrac{a+b+c}{2}$\\
Area: $A=\sqrt{p(p-a)(p-b)(p-c)}$\\
Circumradius: $R=\dfrac{abc}{4A}$\\
Inradius: $r=\dfrac{A}{p}$\\
Length of median (divides triangle into two equal-area triangles): $m_a=\tfrac{1}{2}\sqrt{2b^2+2c^2-a^2}$\\
Length of bisector (divides angles in two): $s_a=\sqrt{bc\left[ 1 - (a / (b+c))^2 \right]}$\\
Law of sines, cosines \& tangents: $\dfrac{\sin\alpha}{a}=\dfrac{\sin\beta}{b}=\dfrac{\sin\gamma}{c}=\dfrac{1}{2R}$\text{.....(1)}\\
$a^2=b^2+c^2-2bc\cos\alpha$\text{.....(2)}\\
$\dfrac{a+b}{a-b}=\dfrac{\tan((\alpha + \beta)/2)}{\tan((\alpha - \beta)/2)}$\text{.....(3)}
\vspace{4mm} % manual %
\subsection{Quadrilaterals}
With side lengths $a,b,c,d$, diagonals $e, f$, diagonals angle $\theta$, area $A$ and
magic flux $F=b^2+d^2-a^2-c^2$:
\[ 4A = 2ef \cdot \sin\theta = F\tan\theta = \sqrt{4e^2f^2-F^2} \]
 For cyclic quadrilaterals the sum of opposite angles is $180^\circ$,
$ef = ac + bd$, and $A = \sqrt{(p-a)(p-b)(p-c)(p-d)}$
\vspace{4mm} % manual %
\subsection{Spherical coordinates}
\begin{center}
\includegraphics[width=25mm]{content/math/sphericalCoordinates}
\end{center}
\[\begin{array}{cc}
x = r\sin\theta\cos\phi & r = \sqrt{x^2+y^2+z^2}\\
y = r\sin\theta\sin\phi & \theta = \textrm{acos}(z/\sqrt{x^2+y^2+z^2})\\
z = r\cos\theta & \phi = \textrm{atan2}(y,x)
\end{array}\]


\section{Sums}
\vspace{-3mm}
\begin{align*}
  c^a + c^{a+1} + \dots + c^{b} &= \frac{c^{b+1} - c^a}{c-1}, \quad c \neq 1\\[1em]
  1^2 + 2^2 + 3^2 + \dots + n^2 &= \frac{n(2n+1)(n+1)}{6}\\
  1^3 + 2^3 + 3^3 + \dots + n^3 &= \frac{n^2(n+1)^2}{4}\\
  1^4 + 2^4 + 3^4 + \dots + n^4 &= \frac{n(n+1)(2n+1)(3n^2+3n-1)}{30}\\[1em]
  1^4 + 2^4 + 3^4 + \dots + n^4 &= S_2 \times \frac{3n^2 + 3n - 1}{5} = S_4
\end{align*}

\small{
  \[b\sum_{k=0}^{n-1}(a + kd)r^k = \frac{ab - (a + nd)br^n}{1-r} + \frac{dbr(1-r^n)}{(1-r)^2}\]\\}
To compute $1^k + ... + n^k$ in $\mathcal{O}(k\lg k + k \lg MOD)$ compute first $t = k + 2$ sums $y_1, ... , y_t$, then interpolate. Let $P = \prod_{i=1}^{t}(n - i)$. Then answer for $n$ is
\[\sum_{i = 1}^{t} \frac{P}{n - i} \cdot \frac{(-1)^{t - i}y_i}{(i - 1)!(t - i)!}\]
Also $S_k = \frac{1}{k+1}\sum_{j = 0}^{k}(-1)^j\binom{k+1}{j}B_j n^{k + 1 - j}$\\
where $B_i$ are Bernoulli numbers.


\section{Bitwise Formulas}

{\small
\subsection*{Some properties of bitwise operations:}
\begin{align*}
  a | b &= a \oplus b + a \& b \\
  a \oplus (a \& b) &= (a | b) \oplus b \\
  b \oplus (a \& b) &= (a | b) \oplus a \\
  (a \& b) \oplus (a | b) &= a \oplus b
\end{align*}

\subsection*{Addition:}
\begin{align*}
  a + b &= a | b + a \& b \\
  a + b &= a \oplus b + 2(a \& b)
\end{align*}

\subsection*{Subtraction:}
\begin{align*}
  a - b &= ((a \& b) \oplus ((a | b) \oplus a)) \\[2mm]
  a - b &= ((a | b) \oplus b) - ((a | b) \oplus a) \\[2mm]
  a - b &= (a \oplus (a \& b)) - ((a | b) \oplus a) \\[2mm]
  a - b &= (a \oplus (a \& b)) - (b \oplus (a \& b)) \\[2mm]
  a - b &= ((a | b) \oplus b) - (b \oplus (a \& b))
\end{align*}

\subsection*{Working rules for logarithms:}
\begin{align*}
  1.\;& \log_b (mn) = \log_b m + \log_b n \\
  2.\;& \log_b \!\left(\frac{m}{n}\right) = \log_b m - \log_b n \\
  3.\;& \log_b (m^r) = r \log_b m \\
  4.\;& \log_b 1 = 0 \\
  5.\;& \log_b b = 1 \\
  6.\;& \log_b m = \frac{\log_a m}{\log_a b}
\end{align*}
}
\section{Trivia}
\textbf{Pythagorean triples}: The Pythagorean triples are uniquely generated by $a=k\cdot (m^{2}-n^{2})$, $b=k\cdot (2mn)$, $c=k\cdot (m^{2}+n^{2})$ with $m > n > 0$, $k > 0$, $\gcd(m,n) = 1$, both $m, n$ not odd.\\
\textbf{Primes}: $p=962592769$ is such that $2^{21} \mid p-1$, which may be useful. For hashing use 970592641 (31-bit number), 31443539979727 (45-bit), 3006703054056749 (52-bit). There are 78498 primes less than 1\,000\,000.\\
\textbf{Primitive roots} modulo $n$ exists iff $n = 1, 2, 4$ or, $n = p^k, 2p^k$ where $p$ is an odd prime. Furthermore, the number of roots are $\phi (\phi (n))$.\\
\textbf{To Find Generator} $g$ of $M$, factor $M - 1$ and get the distinct primes $p_i$. If $g^{(M - 1) / p_i } \neq 1 (MOD M)$ for each $p_i$ then $g$ is a valid root. Try all $g$ until a hit is found (usually found very quick).\\
\textbf{Esitmates}: $\sum_{d|n} d = O(n \log \log n)$.\\
\textbf{Prime count}: 5133 upto 5e4. 9592 upto 1e5. 17984 upto 2e5. 78498 upto 1e6. 5761455 upto 1e8.\\
\textbf{max NOD} $\leq n$: 100 for $n = 5e4$. 500 for $n = 1e7$. 2000 for $n = 1e10$. 200\,000 for $n = 1e19$.\\
\textbf{max Unique Prime Factors}: 6 upto 5e5. 7 upto 9e6. 8 upto 2e8. 9 upto 6e9. 11 upto 7e12. 15 upto 3e19.\\
\textbf{Quadratic Residue}: $(\frac{a}{p})$ is 0 if $p | a$, 1 if $a$ is a quadratic residue, -1 otherwise. Euler: $(\frac{a}{p}) = a^{(p - 1) / 2} (\mod p)$ (prime). Jacobi: if $n = p_1^{e_1}\cdots p_k^{e_k}$ then $(\frac{a}{n}) = \prod (\frac{a}{p_i})^{e_i}$.\\
\textbf{Chicken McNugget.} If $a, b$ coprime, there are $\frac{1}{2}(a - 1)(b - 1)$ numbers not of form $ax + by$ $(x, y \geq 0)$, the largest being $ab - a - b$.

